% Generated by tools/docs/gen_fact_tables.py from rtl/mem/rv32i_clock_gate.sv and docs/manual/common/facts/power.yaml. Do not edit.
{\small
\begin{longtable}{@{}>{\raggedright\arraybackslash}p{31mm}>{\raggedright\arraybackslash}p{9mm}>{\raggedright\arraybackslash}p{19mm}>{\raggedright\arraybackslash}p{22mm}>{\raggedright\arraybackslash}p{58mm}@{}}
\caption{Signals of \code{rv32i_clock_gate}}\label{tab:tm:mod:rv32i-clock-gate:signals}\\
\toprule
Signal & Dir & Width & Clock domain & Description\\
\midrule
\endfirsthead
\multicolumn{5}{@{}l}{\emph{Table \thetable\ continued}}\\
\toprule
Signal & Dir & Width & Clock domain & Description\\
\midrule
\endhead
\bottomrule
\endlastfoot
\code{en} & in & 1 & \code{clk} & Functional enable. The gate takes it into the latch while \code{clk} is low.\\
\code{test_en} & in & 1 & quasi-static & Test enable. High forces the gate open. Keep it static while \code{clk} runs. It can be tied to 0 where no scannable flop sits behind the gate.\\
\code{clk} & in & 1 & clock & Source clock.\\
\code{gclk} & out & 1 & clock & Gated clock.\\
\end{longtable}
}
